\section{Multivariate formal power series}

\begin{definition}
\label{def.fps.mulvar.algebra}
\lean{AlgebraicCombinatorics.FPS, AlgebraicCombinatorics.BivFPS, AlgebraicCombinatorics.BivFPS.x, AlgebraicCombinatorics.BivFPS.y}
\leanhelper
\leanok
Let $k\in\mathbb{N}$. The $K$-algebra of all FPSs in $k$ variables
$x_{1},x_{2},\ldots,x_{k}$ over $K$ will be denoted by $K\left[  \left[
x_{1},x_{2},\ldots,x_{k}\right]  \right]  $.

Sometimes we will use different names for our variables. For example, if we
work with $2$ variables, we will commonly call them $x$ and $y$ instead of
$x_{1}$ and $x_{2}$. Correspondingly, we will use the notation $K\left[
\left[  x,y\right]  \right]  $ for the $K$-algebra of FPSs in these two variables.
\end{definition}

\subsection{Embedding univariate power series into bivariate power series}

\begin{definition}
\label{def.fps.mulvar.embedUnivInBiv}
\lean{AlgebraicCombinatorics.embedUnivInBiv}
\leanhelper
\leanok
Let $f_0, f_1, f_2, \ldots$ be FPSs in a single variable $x$.
Define the bivariate power series
$\operatorname{embed}(f) := \sum_{k \in \mathbb{N}} f_k \, y^k \in K[[x,y]]$,
whose coefficient at $x^n y^k$ is the coefficient of $x^n$ in $f_k$.
\end{definition}

\begin{lemma}
\label{lem.fps.mulvar.coeff-embed}
\lean{AlgebraicCombinatorics.coeff_embedUnivInBiv}
\leanhelper
\leanok
Let $f_0, f_1, f_2, \ldots$ be FPSs in a single variable $x$.
Then for all $n, k \in \mathbb{N}$,
\[
[x^n y^k]\, \operatorname{embed}(f) = [x^n]\, f_k.
\]
\end{lemma}

\begin{proof}
Immediate from the definition of $\operatorname{embed}$.
\end{proof}

\subsection{Comparing coefficients of $y^k$}

\begin{proposition}
\label{prop.fps.mulvar.comp-y-coeff}
\lean{AlgebraicCombinatorics.eq_of_embedUnivInBiv_eq}
\leantarget
\leanok
Let $f_{0},f_{1},f_{2},\ldots$ and
$g_{0},g_{1},g_{2},\ldots$ be FPSs in a single variable $x$ such that
\begin{equation}
\sum_{k\in\mathbb{N}}f_{k}y^{k}=\sum_{k\in\mathbb{N}}g_{k}y^{k}%
\ \ \ \ \ \ \ \ \ \ \text{in }K\left[  \left[  x,y\right]  \right]  .
\label{eq.prop.fps.mulvar.comp-y-coeff.ass}
\end{equation}
Then, $f_{k}=g_{k}$ for each $k\in\mathbb{N}$.
\end{proposition}

\begin{proof}
For each $k\in\mathbb{N}$, let us write the two FPSs $f_{k}$ and $g_{k}$ as
$f_{k}=\sum_{n\in\mathbb{N}}f_{k,n}x^{n}$ and $g_{k}=\sum_{n\in\mathbb{N}%
}g_{k,n}x^{n}$ with $f_{k,n},g_{k,n}\in K$. Then, the equality
(\ref{eq.prop.fps.mulvar.comp-y-coeff.ass}) can be rewritten as
\[
\sum_{k\in\mathbb{N}}\ \ \sum_{n\in\mathbb{N}}f_{k,n}x^{n}y^{k}=\sum
_{k\in\mathbb{N}}\ \ \sum_{n\in\mathbb{N}}g_{k,n}x^{n}y^{k}.
\]
Now, comparing coefficients in front of $x^{n}y^{k}$ in this equality, we
obtain $f_{k,n}=g_{k,n}$ for each $k,n\in\mathbb{N}$. Therefore, $f_{k}=g_{k}$
for each $k\in\mathbb{N}$.
\end{proof}

\subsection{Application: binomial generating function identity}

\begin{definition}
\label{def.fps.mulvar.binomialGenFun}
\lean{AlgebraicCombinatorics.binomialGenFun}
\leanhelper
\leanok
The \emph{binomial generating function} is the bivariate power series
\[
\operatorname{BinGen} := \sum_{n,k\in\mathbb{N}} \binom{n}{k}\, x^n y^k \in \mathbb{Q}[[x,y]],
\]
whose coefficient at $x^n y^k$ is $\binom{n}{k}$.
\end{definition}

\begin{definition}
\label{def.fps.mulvar.binomialGenFunClosedForm}
\lean{AlgebraicCombinatorics.binomialGenFunClosedForm}
\leanhelper
\leanok
The \emph{closed form} of the binomial generating function is
\[
\frac{1}{1 - x(1+y)} \in \mathbb{Q}[[x,y]].
\]
\end{definition}

\begin{theorem}
\label{thm.fps.mulvar.binomialGenFun-eq}
\lean{AlgebraicCombinatorics.binomialGenFun_eq}
\leanhelper
\leanok
The binomial generating function equals its closed form:
\[
\sum_{n,k\in\mathbb{N}} \binom{n}{k}\, x^n y^k = \frac{1}{1-x(1+y)}.
\]
\end{theorem}

\begin{proof}
\leanok
It suffices to show that $(1 - x(1+y)) \cdot \operatorname{BinGen} = 1$, since
the constant term of $1 - x(1+y)$ is $1 \neq 0$ and inverses in power series
rings are unique.

Since $1 - x(1+y) = 1 - x - xy$, the coefficient of $x^n y^k$ in
$(1 - x - xy) \cdot \operatorname{BinGen}$ is
\[
\binom{n}{k} - \binom{n-1}{k} - \binom{n-1}{k-1}
\]
(where we set $\binom{n-1}{k} = 0$ when $n = 0$ and $\binom{n-1}{k-1} = 0$
when $n = 0$ or $k = 0$).

For $n = 0$, this equals $\binom{0}{k} = \delta_{k,0}$, which is the
coefficient of $x^0 y^k$ in $1$.

For $n \geq 1$ and $k = 0$, this equals $\binom{n}{0} - \binom{n-1}{0} = 1 - 1 = 0$.

For $n \geq 1$ and $k \geq 1$, Pascal's identity gives
$\binom{n}{k} = \binom{n-1}{k-1} + \binom{n-1}{k}$, so the expression is $0$.

Thus $(1 - x(1+y)) \cdot \operatorname{BinGen} = 1$, and uniqueness of inverses
yields $\operatorname{BinGen} = (1 - x(1+y))^{-1}$.
\end{proof}

From the computation in the source text, comparing the two expressions for
$\sum_{n,k} \binom{n}{k} x^n y^k$, we find
\begin{equation}
\sum_{k\in\mathbb{N}}\dfrac{x^{k}}{\left(  1-x\right)  ^{k+1}}y^{k}=\sum
_{k\in\mathbb{N}}\left(  \sum_{n\in\mathbb{N}}\dbinom{n}{k}x^{n}\right)
y^{k}. \label{eq.fps.mulvar.exa1.xyk}
\end{equation}

\subsection{Helper lemmas from the formalization}

\begin{lemma}
\label{lem.fps.mulvar.one-sub-X-inv}
\lean{AlgebraicCombinatorics.one_sub_X_inv_eq_mk_one}
\leanhelper
\leanok
The inverse of $1 - x$ in $\mathbb{Q}[[x]]$ equals the power series with all
coefficients equal to $1$:
\[
(1 - x)^{-1} = \sum_{n \in \mathbb{N}} x^n.
\]
\end{lemma}

\begin{proof}
We verify $(1-x) \cdot \sum_{n \in \mathbb{N}} x^n = 1$ by direct
computation, and the constant term of $1-x$ is $1 \neq 0$.
\end{proof}

\begin{lemma}
\label{lem.fps.mulvar.X-pow-mul-inv}
\lean{AlgebraicCombinatorics.X_pow_mul_inv_one_sub_pow_eq_mk_choose}
\leanhelper
\leanok
For each $k \in \mathbb{N}$,
\[
\frac{x^k}{(1-x)^{k+1}} = \sum_{n \in \mathbb{N}} \binom{n}{k}\, x^n
\quad \text{in } \mathbb{Q}[[x]].
\]
\end{lemma}

\begin{proof}
We compare coefficients. By the formula for the inverse of $(1-x)^{k+1}$,
$(1-x)^{-(k+1)} = \sum_n \binom{n+k}{k}\, x^n$.
Multiplying by $x^k$ shifts the index: the coefficient of $x^n$ in
$x^k \cdot (1-x)^{-(k+1)}$ is $0$ for $n < k$ and $\binom{n-k+k}{k} = \binom{n}{k}$
for $n \geq k$. For $n < k$, $\binom{n}{k} = 0$ as well, so both sides agree.
\end{proof}

\begin{theorem}
\label{thm.fps.mulvar.exa1.xyk}
\lean{AlgebraicCombinatorics.binomialGenFun_xyk_eq}
\leanhelper
\leanok
Equation~(\ref{eq.fps.mulvar.exa1.xyk}) holds: as embedded sequences of
univariate power series,
\[
\operatorname{embed}\!\left(k \mapsto \frac{x^k}{(1-x)^{k+1}}\right)
= \operatorname{embed}\!\left(k \mapsto \sum_{n \in \mathbb{N}} \binom{n}{k}\, x^n\right).
\]
\end{theorem}

\begin{proof}
Follows immediately from Lemma~\ref{lem.fps.mulvar.X-pow-mul-inv}, which
shows that the two sequences are equal term by term.
\end{proof}

\begin{theorem}
\label{thm.fps.mulvar.exa1.res1}
\lean{AlgebraicCombinatorics.sum_choose_pow_eq}
\leanhelper
\leanok
For each $k \in \mathbb{N}$,
\[
\dfrac{x^{k}}{\left(  1-x\right)  ^{k+1}}=\sum_{n\in\mathbb{N}}\dbinom{n}%
{k}x^{n}
\]
as an identity in $\mathbb{Q}[[x]]$.
This is an equality between \textbf{univariate} FPSs, even though we have
obtained it by manipulating \textbf{bivariate} FPSs.
\end{theorem}

\begin{proof}
By Theorem~\ref{thm.fps.mulvar.exa1.xyk}, the two embedded sequences are
equal. Apply Proposition~\ref{prop.fps.mulvar.comp-y-coeff}
to extract the $k$-th component of
this equality.
\end{proof}

\subsection{Partial derivatives}

\begin{definition}
\label{def.fps.mulvar.partialDeriv}
\lean{AlgebraicCombinatorics.partialDeriv}
\leanhelper
\leanok
Let $\sigma$ be a type of variable indices and let $i \in \sigma$.
The \emph{partial derivative} of a multivariate power series
$f = \sum_{\mathbf{m}} a_{\mathbf{m}}\, \mathbf{x}^{\mathbf{m}}$
with respect to $x_i$ is the power series
\[
\frac{\partial f}{\partial x_i}
= \sum_{\mathbf{m}} (m_i + 1) \cdot a_{\mathbf{m} + \mathbf{e}_i}
  \cdot \mathbf{x}^{\mathbf{m}},
\]
where $\mathbf{e}_i$ is the $i$-th standard basis vector
(i.e., $\mathbf{e}_i = (\delta_{j,i})_{j \in \sigma}$).

This defines an $R$-linear map $\partial/\partial x_i \colon
R[[\mathbf{x}]] \to R[[\mathbf{x}]]$.
\end{definition}

\begin{lemma}
\label{lem.fps.mulvar.coeff-partialDeriv}
\lean{AlgebraicCombinatorics.coeff_partialDeriv}
\leanhelper
\leanok
For any multivariate power series $f$ and any multi-index $\mathbf{m}$,
\[
[\mathbf{x}^{\mathbf{m}}]\, \frac{\partial f}{\partial x_i}
= (m_i + 1) \cdot [\mathbf{x}^{\mathbf{m} + \mathbf{e}_i}]\, f.
\]
\end{lemma}

\begin{proof}
Immediate from the definition.
\end{proof}

\begin{lemma}
\label{lem.fps.mulvar.antidiag-shift-fst}
\lean{AlgebraicCombinatorics.antidiag_shift_fst}
\leanhelper
\leanok
The map $(\mathbf{a}, \mathbf{b}) \mapsto (\mathbf{a} + \mathbf{e}_i, \mathbf{b})$ gives a bijection from
$\operatorname{antidiagonal}(\mathbf{m})$ to the subset of
$\operatorname{antidiagonal}(\mathbf{m} + \mathbf{e}_i)$ consisting of
pairs $(\mathbf{a}, \mathbf{b})$ with $\mathbf{e}_i \leq \mathbf{a}$.
\end{lemma}

\begin{proof}
Verified by checking the membership conditions in both directions.
\end{proof}

\begin{lemma}
\label{lem.fps.mulvar.antidiag-shift-snd}
\lean{AlgebraicCombinatorics.antidiag_shift_snd}
\leanhelper
\leanok
The map $(\mathbf{a}, \mathbf{b}) \mapsto (\mathbf{a}, \mathbf{b} + \mathbf{e}_i)$ gives a bijection from
$\operatorname{antidiagonal}(\mathbf{m})$ to the subset of
$\operatorname{antidiagonal}(\mathbf{m} + \mathbf{e}_i)$ consisting of
pairs $(\mathbf{a}, \mathbf{b})$ with $\mathbf{e}_i \leq \mathbf{b}$.
\end{lemma}

\begin{proof}
Symmetric to Lemma~\ref{lem.fps.mulvar.antidiag-shift-fst}.
\end{proof}

\begin{lemma}
\label{lem.fps.mulvar.sum-filter-zero-fst}
\lean{AlgebraicCombinatorics.sum_filter_zero_fst}
\leanhelper
\leanok
In the sum over $\operatorname{antidiagonal}(\mathbf{m} + \mathbf{e}_i)$,
the terms with $a_i = 0$ (i.e., $\mathbf{e}_i \not\leq \mathbf{a}$) contribute zero
to $\sum_{(\mathbf{a},\mathbf{b})} a_i \cdot [{\mathbf{x}}^{\mathbf{a}}]\, f \cdot [{\mathbf{x}}^{\mathbf{b}}]\, g$.
\end{lemma}

\begin{proof}
Each such term has $a_i = 0$, so the factor $a_i$ makes the whole
summand vanish.
\end{proof}

\begin{lemma}
\label{lem.fps.mulvar.sum-filter-zero-snd}
\lean{AlgebraicCombinatorics.sum_filter_zero_snd}
\leanhelper
\leanok
In the sum over $\operatorname{antidiagonal}(\mathbf{m} + \mathbf{e}_i)$,
the terms with $b_i = 0$ (i.e., $\mathbf{e}_i \not\leq \mathbf{b}$) contribute zero
to $\sum_{(\mathbf{a},\mathbf{b})} [\mathbf{x}^{\mathbf{a}}]\, f \cdot b_i \cdot [\mathbf{x}^{\mathbf{b}}]\, g$.
\end{lemma}

\begin{proof}
Each such term has $b_i = 0$, so the factor $b_i$ makes the whole
summand vanish.
\end{proof}

\begin{theorem}[Product rule for partial derivatives]
\label{thm.fps.mulvar.partialDeriv-mul}
\lean{AlgebraicCombinatorics.partialDeriv_mul}
\leanhelper
\leanok
The partial derivative satisfies the product rule:
\[
\frac{\partial (fg)}{\partial x_i}
= \frac{\partial f}{\partial x_i} \cdot g
  + f \cdot \frac{\partial g}{\partial x_i}.
\]
\end{theorem}

\begin{proof}
We compare coefficients at an arbitrary multi-index $\mathbf{m}$.
Set $\mathbf{m}' = \mathbf{m} + \mathbf{e}_i$.
On the left side, by Lemma~\ref{lem.fps.mulvar.coeff-partialDeriv} and the
product formula,
\[
[\mathbf{x}^{\mathbf{m}}]\, \frac{\partial(fg)}{\partial x_i}
= m'_i \cdot \sum_{\substack{\mathbf{a}+\mathbf{b}=\mathbf{m}'}}
  [\mathbf{x}^{\mathbf{a}}]\, f \cdot [\mathbf{x}^{\mathbf{b}}]\, g.
\]
We distribute the scalar $m'_i$ into each summand. For each pair
$(\mathbf{a},\mathbf{b})$ in $\operatorname{antidiagonal}(\mathbf{m}')$,
we have $m'_i = a_i + b_i$, so
\[
m'_i \cdot [\mathbf{x}^{\mathbf{a}}]\, f \cdot [\mathbf{x}^{\mathbf{b}}]\, g
= a_i \cdot [\mathbf{x}^{\mathbf{a}}]\, f \cdot [\mathbf{x}^{\mathbf{b}}]\, g
  + [\mathbf{x}^{\mathbf{a}}]\, f \cdot b_i \cdot [\mathbf{x}^{\mathbf{b}}]\, g.
\]
Summing over the antidiagonal, the first sum equals
$[\mathbf{x}^{\mathbf{m}}]\, (\frac{\partial f}{\partial x_i} \cdot g)$
and the second equals
$[\mathbf{x}^{\mathbf{m}}]\, (f \cdot \frac{\partial g}{\partial x_i})$,
after reindexing via Lemmas~\ref{lem.fps.mulvar.antidiag-shift-fst}
and~\ref{lem.fps.mulvar.antidiag-shift-snd} and discarding the zero
terms via Lemmas~\ref{lem.fps.mulvar.sum-filter-zero-fst}
and~\ref{lem.fps.mulvar.sum-filter-zero-snd}.
\end{proof}
